% Options for packages loaded elsewhere
\PassOptionsToPackage{unicode}{hyperref}
\PassOptionsToPackage{hyphens}{url}
\documentclass[
]{article}
\usepackage{xcolor}
\usepackage[margin=1in]{geometry}
\usepackage{amsmath,amssymb}
\setcounter{secnumdepth}{-\maxdimen} % remove section numbering
\usepackage{iftex}
\ifPDFTeX
  \usepackage[T1]{fontenc}
  \usepackage[utf8]{inputenc}
  \usepackage{textcomp} % provide euro and other symbols
\else % if luatex or xetex
  \usepackage{unicode-math} % this also loads fontspec
  \defaultfontfeatures{Scale=MatchLowercase}
  \defaultfontfeatures[\rmfamily]{Ligatures=TeX,Scale=1}
\fi
\usepackage{lmodern}
\ifPDFTeX\else
  % xetex/luatex font selection
\fi
% Use upquote if available, for straight quotes in verbatim environments
\IfFileExists{upquote.sty}{\usepackage{upquote}}{}
\IfFileExists{microtype.sty}{% use microtype if available
  \usepackage[]{microtype}
  \UseMicrotypeSet[protrusion]{basicmath} % disable protrusion for tt fonts
}{}
\makeatletter
\@ifundefined{KOMAClassName}{% if non-KOMA class
  \IfFileExists{parskip.sty}{%
    \usepackage{parskip}
  }{% else
    \setlength{\parindent}{0pt}
    \setlength{\parskip}{6pt plus 2pt minus 1pt}}
}{% if KOMA class
  \KOMAoptions{parskip=half}}
\makeatother
\usepackage{graphicx}
\makeatletter
\newsavebox\pandoc@box
\newcommand*\pandocbounded[1]{% scales image to fit in text height/width
  \sbox\pandoc@box{#1}%
  \Gscale@div\@tempa{\textheight}{\dimexpr\ht\pandoc@box+\dp\pandoc@box\relax}%
  \Gscale@div\@tempb{\linewidth}{\wd\pandoc@box}%
  \ifdim\@tempb\p@<\@tempa\p@\let\@tempa\@tempb\fi% select the smaller of both
  \ifdim\@tempa\p@<\p@\scalebox{\@tempa}{\usebox\pandoc@box}%
  \else\usebox{\pandoc@box}%
  \fi%
}
% Set default figure placement to htbp
\def\fps@figure{htbp}
\makeatother
\setlength{\emergencystretch}{3em} % prevent overfull lines
\providecommand{\tightlist}{%
  \setlength{\itemsep}{0pt}\setlength{\parskip}{0pt}}
\usepackage{bookmark}
\IfFileExists{xurl.sty}{\usepackage{xurl}}{} % add URL line breaks if available
\urlstyle{same}
\hypersetup{
  hidelinks,
  pdfcreator={LaTeX via pandoc}}

\author{}
\date{\vspace{-2.5em}}

\begin{document}

\section{📱 TEMPLATE DE PUBLICACIÓN LINKEDIN - DICIEMBRE
2025}\label{template-de-publicaciuxf3n-linkedin---diciembre-2025}

\subsection{🎯 ESTRATEGIA BASADA EN MEJORES PRÁCTICAS
2025}\label{estrategia-basada-en-mejores-pruxe1cticas-2025}

\subsubsection{Principios Clave (Actualizados Diciembre
2025)}\label{principios-clave-actualizados-diciembre-2025}

\begin{enumerate}
\def\labelenumi{\arabic{enumi}.}
\tightlist
\item
  ✅ \textbf{NO incluir enlaces en el post principal} (penaliza el
  algoritmo)
\item
  ✅ \textbf{Enlaces en el PRIMER COMENTARIO} (mejor alcance)
\item
  ✅ \textbf{Formato vertical} para videos/imágenes (mayor engagement)
\item
  ✅ \textbf{Carruseles} tienen mejor rendimiento que posts simples
\item
  ✅ \textbf{Storytelling} \textgreater{} Promoción directa
\item
  ✅ \textbf{Call-to-action claro} pero no agresivo
\item
  ✅ \textbf{Hashtags estratégicos} (3-5 máximo, relevantes)
\end{enumerate}

\begin{center}\rule{0.5\linewidth}{0.5pt}\end{center}

\subsection{📝 VERSIÓN 1: POST PRINCIPAL
(STORYTELLING)}\label{versiuxf3n-1-post-principal-storytelling}

\subsubsection{Texto del Post}\label{texto-del-post}

\begin{verbatim}
🎯 ¿Y si te dijera que puedes generar 300+ versiones ÚNICAS del mismo ejercicio matemático en segundos?

Hace meses comencé un proyecto que parecía imposible:
Crear un sistema que revolucionara la forma de evaluar matemáticas en Colombia.

El problema era claro:
❌ Profesores creando ejercicios manualmente (horas de trabajo)
❌ Estudiantes memorizando respuestas en lugar de entender conceptos
❌ Evaluaciones estáticas que no se adaptan a cada alumno

La solución: R/exams + Python + TikZ

Hoy comparto el primer demo público del sistema:
📊 "Consumo Telefónico Adicional" - Nivel ICFES 2

Lo que hace especial a este ejercicio:

✅ 300+ versiones completamente diferentes
✅ Gráficos dinámicos generados con Python
✅ Verificación automática de respuestas
✅ Explicaciones paso a paso personalizadas
✅ Compatible con Moodle, PDF, HTML interactivo

Cada versión tiene:
→ Datos aleatorios pero coherentes matemáticamente
→ Gráficos únicos generados automáticamente
→ Distractores pedagógicos (no respuestas al azar)
→ Contextos realistas y relevantes

¿El resultado?
Estudiantes que realmente ENTIENDEN en lugar de memorizar.
Profesores que ahorran 90% del tiempo en crear evaluaciones.

He preparado 5 demos interactivos donde puedes:
🔹 Resolver el ejercicio directamente en tu navegador
🔹 Verificar tu respuesta al instante
🔹 Ver la solución completa explicada
🔹 Probar diferentes versiones del mismo concepto

También incluyo:
📄 PDF con 10 versiones para imprimir
💻 Código fuente completo en GitHub
🎓 Archivo Moodle XML listo para importar

Este es solo el PRIMERO de muchos ejercicios que compartiré.

¿Te interesa ver cómo funciona?
👇 Todos los enlaces en el primer comentario

#EducaciónMatemática #ICFES #InnovaciónEducativa #RStats #Python
\end{verbatim}

\begin{center}\rule{0.5\linewidth}{0.5pt}\end{center}

\subsection{💬 PRIMER COMENTARIO (CON
ENLACES)}\label{primer-comentario-con-enlaces}

\subsubsection{Texto del Comentario}\label{texto-del-comentario}

\begin{verbatim}
🔗 ACCEDE A LOS DEMOS INTERACTIVOS:

📊 Página principal con todos los demos:
https://alvaretto.github.io/proyecto-r-exams-icfes-matematicas-optimizado/

🎯 Demos individuales (elige tu versión favorita):

Versión 1: https://alvaretto.github.io/proyecto-r-exams-icfes-matematicas-optimizado/demos/demo_consumo_telefonico_v1.html

Versión 2: https://alvaretto.github.io/proyecto-r-exams-icfes-matematicas-optimizado/demos/demo_consumo_telefonico_v2.html

Versión 3: https://alvaretto.github.io/proyecto-r-exams-icfes-matematicas-optimizado/demos/demo_consumo_telefonico_v3.html

Versión 4: https://alvaretto.github.io/proyecto-r-exams-icfes-matematicas-optimizado/demos/demo_consumo_telefonico_v4.html

Versión 5: https://alvaretto.github.io/proyecto-r-exams-icfes-matematicas-optimizado/demos/demo_consumo_telefonico_v5.html

📥 RECURSOS DESCARGABLES:

📄 PDF con 10 versiones:
https://alvaretto.github.io/proyecto-r-exams-icfes-matematicas-optimizado/recursos/muestra_10_versiones_consumo_telefonico.pdf

🎓 Archivo Moodle XML:
https://alvaretto.github.io/proyecto-r-exams-icfes-matematicas-optimizado/recursos/consumo_telefonico_moodle.xml

💻 Código fuente completo (.Rmd):
https://github.com/alvaretto/proyecto-r-exams-icfes-matematicas-optimizado

---

💡 TIP: Prueba resolver el mismo ejercicio en 2-3 versiones diferentes.
Notarás que aunque los números cambian, el concepto matemático es el mismo.
Así es como se aprende de verdad.

¿Preguntas? ¡Déjalas en los comentarios! 👇
\end{verbatim}

\begin{center}\rule{0.5\linewidth}{0.5pt}\end{center}

\subsection{🎨 VERSIÓN 2: POST CON ENFOQUE
TÉCNICO}\label{versiuxf3n-2-post-con-enfoque-tuxe9cnico}

\subsubsection{Para Audiencia de Desarrolladores/Educadores
Técnicos}\label{para-audiencia-de-desarrolladoreseducadores-tuxe9cnicos}

\begin{verbatim}
🚀 OPEN SOURCE: Sistema de Evaluación Matemática con R/exams

Stack tecnológico:
→ R/exams (generación dinámica)
→ Python + matplotlib (gráficos)
→ TikZ + LaTeX (diagramas profesionales)
→ exams2forms (demos interactivos HTML)
→ GitHub Pages (hosting gratuito)

Características del sistema:

📊 ALEATORIZACIÓN INTELIGENTE
• 300+ versiones únicas verificadas con tests automáticos
• Datos coherentes matemáticamente (no aleatorización superficial)
• Distractores basados en errores conceptuales reales

🎯 MULTI-FORMATO
• HTML interactivo (verificación automática)
• PDF para impresión
• Moodle XML (importación directa)
• NOPS (exámenes escritos con lectura óptica)

🔧 CÓDIGO REPRODUCIBLE
• 100% código abierto en GitHub
• Documentación completa
• Tests automatizados
• Ejemplos funcionales

📈 MÉTRICAS DE CALIDAD
• Fidelidad visual 98%+ en gráficos TikZ
• Diversidad de versiones validada automáticamente
• Metadatos ICFES completos
• Tolerancias configurables para evaluación automática

Primer demo público disponible:
"Consumo Telefónico Adicional" - Interpretación de gráficos + cálculos

Incluye:
✅ 5 demos HTML interactivos
✅ PDF con 10 versiones
✅ Código fuente .Rmd completo
✅ Archivo Moodle XML

Enlaces en el primer comentario 👇

#RStats #Python #OpenSource #EdTech #DataScience
\end{verbatim}

\begin{center}\rule{0.5\linewidth}{0.5pt}\end{center}

\subsection{📸 ELEMENTOS VISUALES
RECOMENDADOS}\label{elementos-visuales-recomendados}

\subsubsection{Opción 1: Carrusel de Imágenes
(RECOMENDADO)}\label{opciuxf3n-1-carrusel-de-imuxe1genes-recomendado}

\textbf{Slide 1}: Captura del demo interactivo mostrando el gráfico de
barras \textbf{Slide 2}: Captura mostrando la tabla de estado de cuenta
\textbf{Slide 3}: Captura del botón de verificación con feedback
``Correcto'' \textbf{Slide 4}: Captura de la solución completa explicada
\textbf{Slide 5}: Captura del código .Rmd (primeras líneas)
\textbf{Slide 6}: Infografía con estadísticas (300+ versiones, 5
formatos, etc.)

\subsubsection{Opción 2: Video Corto (60-90
segundos)}\label{opciuxf3n-2-video-corto-60-90-segundos}

\textbf{Guion sugerido}:

\begin{verbatim}
0:00-0:10 → Mostrar problema: "Profesores creando ejercicios manualmente"
0:10-0:20 → Introducir solución: "Sistema R/exams"
0:20-0:40 → Demo en vivo: Resolver ejercicio, verificar respuesta
0:40-0:50 → Mostrar diferentes versiones del mismo ejercicio
0:50-0:60 → Mostrar código fuente brevemente
0:60-0:90 → Call-to-action: "Enlaces en comentarios"
\end{verbatim}

\subsubsection{Opción 3: Imagen Única con
Infografía}\label{opciuxf3n-3-imagen-uxfanica-con-infografuxeda}

\textbf{Elementos a incluir}:

\begin{itemize}
\tightlist
\item
  Logo del proyecto
\item
  Título: ``Sistema R/exams para Matemáticas ICFES''
\item
  Iconos con estadísticas clave:

  \begin{itemize}
  \tightlist
  \item
    300+ versiones únicas
  \item
    5 formatos de salida
  \item
    100\% código abierto
  \item
    Compatible con Moodle
  \end{itemize}
\item
  Screenshot del demo interactivo
\item
  QR code a la página principal (opcional)
\end{itemize}

\begin{center}\rule{0.5\linewidth}{0.5pt}\end{center}

\subsection{🎯 HASHTAGS ESTRATÉGICOS}\label{hashtags-estratuxe9gicos}

\subsubsection{Hashtags Principales (Usar
3-5)}\label{hashtags-principales-usar-3-5}

\textbf{Para Educadores}:

\begin{verbatim}
#EducaciónMatemática #ICFES #InnovaciónEducativa #EdTech #Moodle
\end{verbatim}

\textbf{Para Desarrolladores}:

\begin{verbatim}
#RStats #Python #OpenSource #DataScience #GitHub
\end{verbatim}

\textbf{Para Instituciones}:

\begin{verbatim}
#EducaciónSuperior #TransformaciónDigital #AprendizajeAdaptativo
\end{verbatim}

\subsubsection{Hashtags Secundarios
(Rotar)}\label{hashtags-secundarios-rotar}

\begin{verbatim}
#Estadística #Matemáticas #ProgramaciónR #DataVisualization
#EvaluaciónEducativa #TecnologíaEducativa #CódigoAbierto
#EducaciónColombia #PreICFES #DocentesInnovadores
\end{verbatim}

\begin{center}\rule{0.5\linewidth}{0.5pt}\end{center}

\subsection{📅 CALENDARIO DE
PUBLICACIÓN}\label{calendario-de-publicaciuxf3n}

\subsubsection{Timing Óptimo (Basado en Datos
2025)}\label{timing-uxf3ptimo-basado-en-datos-2025}

\textbf{Mejores días}: Martes, Miércoles, Jueves \textbf{Mejores
horarios}:

\begin{itemize}
\tightlist
\item
  🌅 7:00-9:00 AM (antes de iniciar jornada laboral)
\item
  🌆 12:00-1:00 PM (hora de almuerzo)
\item
  🌃 5:00-6:00 PM (fin de jornada laboral)
\end{itemize}

\textbf{Recomendación para primera publicación}: 📅 \textbf{Martes o
Miércoles a las 8:00 AM} (hora Colombia)

\begin{center}\rule{0.5\linewidth}{0.5pt}\end{center}

\subsection{🎬 PLAN DE LANZAMIENTO (PRIMERA
SEMANA)}\label{plan-de-lanzamiento-primera-semana}

\subsubsection{Día 1 (Martes 8:00 AM): Publicación
Principal}\label{duxeda-1-martes-800-am-publicaciuxf3n-principal}

\begin{itemize}
\tightlist
\item
  ✅ Post con storytelling (Versión 1)
\item
  ✅ Carrusel de 6 imágenes
\item
  ✅ Primer comentario con todos los enlaces
\item
  ✅ Responder TODOS los comentarios en las primeras 2 horas
\end{itemize}

\subsubsection{Día 2 (Miércoles 12:00 PM): Engagement
Boost}\label{duxeda-2-miuxe9rcoles-1200-pm-engagement-boost}

\begin{itemize}
\item
  ✅ Comentar en el post original con estadística interesante:

\begin{verbatim}
📊 UPDATE: En 24 horas, más de [X] personas han probado los demos.

El feedback más común: "No sabía que esto era posible con R"

¿Qué te sorprendió más del sistema? 👇
\end{verbatim}
\end{itemize}

\subsubsection{Día 3 (Jueves 8:00 AM): Post de
Seguimiento}\label{duxeda-3-jueves-800-am-post-de-seguimiento}

\begin{itemize}
\tightlist
\item
  ✅ Post técnico (Versión 2) para audiencia de desarrolladores
\item
  ✅ Enlazar al post original
\item
  ✅ Compartir métricas de engagement del primer post
\end{itemize}

\subsubsection{Día 5 (Sábado 10:00 AM): Behind the
Scenes}\label{duxeda-5-suxe1bado-1000-am-behind-the-scenes}

\begin{itemize}
\tightlist
\item
  ✅ Post mostrando el proceso de creación
\item
  ✅ Código fuente comentado
\item
  ✅ Desafíos técnicos superados
\end{itemize}

\begin{center}\rule{0.5\linewidth}{0.5pt}\end{center}

\subsection{💡 ESTRATEGIAS DE
ENGAGEMENT}\label{estrategias-de-engagement}

\subsubsection{Preguntas para Generar
Conversación}\label{preguntas-para-generar-conversaciuxf3n}

En el post principal o comentarios:

\begin{verbatim}
❓ ¿Cuánto tiempo te toma crear un ejercicio matemático desde cero?

❓ ¿Qué herramientas usas actualmente para evaluar a tus estudiantes?

❓ ¿Conocías R/exams antes de este post?

❓ ¿Qué tipo de ejercicios te gustaría ver próximamente?
   (Geometría, Álgebra, Estadística, Cálculo...)

❓ ¿Usas Moodle u otro LMS en tu institución?
\end{verbatim}

\subsubsection{Respuestas Preparadas para Comentarios
Comunes}\label{respuestas-preparadas-para-comentarios-comunes}

\textbf{``¿Esto funciona con otros temas además de matemáticas?''}

\begin{verbatim}
¡Excelente pregunta! R/exams es agnóstico al contenido.
Funciona para cualquier materia que pueda evaluarse con:
• Opción múltiple
• Respuestas numéricas
• Respuestas de texto
• Combinaciones (Cloze)

He visto implementaciones en: Física, Química, Economía,
Estadística, Programación, incluso Idiomas.

¿En qué materia te gustaría implementarlo?
\end{verbatim}

\textbf{``¿Es difícil de aprender?''}

\begin{verbatim}
Curva de aprendizaje moderada si ya conoces R.

Si eres nuevo en R:
→ 2-3 semanas para ejercicios básicos
→ 1-2 meses para ejercicios avanzados con gráficos

Ventaja: Una vez aprendes, puedes generar miles de
ejercicios con mínimo esfuerzo adicional.

Compartiré tutoriales paso a paso próximamente.
¿Te interesaría?
\end{verbatim}

\textbf{``¿Dónde puedo aprender más?''}

\begin{verbatim}
Recursos que recomiendo:

📚 Documentación oficial: https://www.r-exams.org/
💻 Mi repositorio GitHub: [enlace]
📹 Tutoriales en video: [si tienes]

También planeo compartir:
→ Ejercicios nuevos semanalmente
→ Tutoriales específicos
→ Casos de uso reales

¿Qué te gustaría aprender primero?
\end{verbatim}

\begin{center}\rule{0.5\linewidth}{0.5pt}\end{center}

\subsection{📊 MÉTRICAS DE ÉXITO A
MONITOREAR}\label{muxe9tricas-de-uxe9xito-a-monitorear}

\subsubsection{KPIs Principales (Primeras 48
horas)}\label{kpis-principales-primeras-48-horas}

\begin{itemize}
\tightlist
\item
  🎯 \textbf{Impresiones}: Meta 5,000+
\item
  👁️ \textbf{Visualizaciones}: Meta 1,000+
\item
  💬 \textbf{Comentarios}: Meta 20+
\item
  🔄 \textbf{Compartidos}: Meta 10+
\item
  👍 \textbf{Reacciones}: Meta 100+
\item
  🔗 \textbf{Clicks en enlaces}: Meta 50+
\end{itemize}

\subsubsection{KPIs Secundarios (Primera
semana)}\label{kpis-secundarios-primera-semana}

\begin{itemize}
\tightlist
\item
  📈 \textbf{Nuevos seguidores}: Meta 50+
\item
  🌐 \textbf{Visitas a GitHub Pages}: Meta 200+
\item
  ⭐ \textbf{Stars en GitHub}: Meta 10+
\item
  📥 \textbf{Descargas de recursos}: Meta 30+
\end{itemize}

\subsubsection{Herramientas de
Seguimiento}\label{herramientas-de-seguimiento}

\begin{verbatim}
✅ LinkedIn Analytics (nativo)
✅ Google Analytics en GitHub Pages
✅ GitHub Insights (traffic, clones, stars)
✅ Bitly (para acortar y trackear enlaces - opcional)
\end{verbatim}

\begin{center}\rule{0.5\linewidth}{0.5pt}\end{center}

\subsection{🎨 PLANTILLA DE IMÁGENES PARA
CARRUSEL}\label{plantilla-de-imuxe1genes-para-carrusel}

\subsubsection{Especificaciones
Técnicas}\label{especificaciones-tuxe9cnicas}

\begin{itemize}
\tightlist
\item
  \textbf{Formato}: PNG o JPG
\item
  \textbf{Dimensiones}: 1080x1080 px (cuadrado) o 1080x1350 px
  (vertical)
\item
  \textbf{Peso máximo}: 10 MB por imagen
\item
  \textbf{Cantidad}: 6-10 slides (óptimo: 6)
\end{itemize}

\subsubsection{Contenido Sugerido por
Slide}\label{contenido-sugerido-por-slide}

\textbf{Slide 1 - Portada}:

\begin{verbatim}
Título: "Sistema R/exams para Matemáticas ICFES"
Subtítulo: "300+ versiones únicas del mismo ejercicio"
Imagen de fondo: Screenshot del demo interactivo
Logo del proyecto
\end{verbatim}

\textbf{Slide 2 - El Problema}:

\begin{verbatim}
Título: "El Problema"
Bullet points:
• Profesores creando ejercicios manualmente
• Estudiantes memorizando en lugar de entender
• Evaluaciones estáticas y repetitivas
Icono: ❌ o 😓
\end{verbatim}

\textbf{Slide 3 - La Solución}:

\begin{verbatim}
Título: "La Solución"
Stack tecnológico con iconos:
• R/exams
• Python + matplotlib
• TikZ + LaTeX
• GitHub Pages
Icono: ✅ o 🚀
\end{verbatim}

\textbf{Slide 4 - Demo en Acción}:

\begin{verbatim}
Screenshot del demo interactivo mostrando:
• Gráfico de barras
• Tabla de datos
• Botones de interacción
Anotaciones con flechas señalando características clave
\end{verbatim}

\textbf{Slide 5 - Características Clave}:

\begin{verbatim}
Título: "¿Qué lo hace especial?"
Iconos con texto:
✅ 300+ versiones únicas
✅ Verificación automática
✅ Explicaciones paso a paso
✅ Compatible con Moodle
✅ 100% código abierto
\end{verbatim}

\textbf{Slide 6 - Call to Action}:

\begin{verbatim}
Título: "Pruébalo Ahora"
Texto: "Enlaces en el primer comentario 👇"
QR code a la página principal (opcional)
Iconos de redes sociales
\end{verbatim}

\begin{center}\rule{0.5\linewidth}{0.5pt}\end{center}

\subsection{🔄 ESTRATEGIA DE REUTILIZACIÓN DE
CONTENIDO}\label{estrategia-de-reutilizaciuxf3n-de-contenido}

\subsubsection{Adaptaciones para Otras
Plataformas}\label{adaptaciones-para-otras-plataformas}

\textbf{Twitter/X}:

\begin{verbatim}
🎯 Nuevo proyecto: Sistema R/exams para generar 300+ versiones
únicas de ejercicios matemáticos ICFES

✅ Verificación automática
✅ Gráficos dinámicos con Python
✅ Compatible con Moodle
✅ 100% código abierto

Demos interactivos: [enlace acortado]

#RStats #Python #EdTech
\end{verbatim}

\textbf{Facebook (Grupos de Educadores)}:

\begin{verbatim}
[Versión más extensa del post de LinkedIn]
[Incluir enlaces directamente en el post]
[Agregar más contexto educativo]
\end{verbatim}

\textbf{Instagram (Stories)}:

\begin{verbatim}
Story 1: Screenshot del demo + "Swipe up para probar"
Story 2: Video corto resolviendo el ejercicio
Story 3: Código fuente con música de fondo
Story 4: Estadísticas (300+ versiones, etc.)
Story 5: Call to action con enlace
\end{verbatim}

\textbf{YouTube (Video Tutorial)}:

\begin{verbatim}
Título: "Cómo Generar 300+ Ejercicios Matemáticos Únicos con R/exams"
Duración: 10-15 minutos
Contenido:

- Introducción al problema (2 min)
- Demo del sistema (3 min)
- Explicación técnica (5 min)
- Cómo empezar (3 min)
- Call to action (1 min)
\end{verbatim}

\begin{center}\rule{0.5\linewidth}{0.5pt}\end{center}

\subsection{📧 EMAIL MARKETING
(OPCIONAL)}\label{email-marketing-opcional}

\subsubsection{Para Lista de Correo de
Educadores}\label{para-lista-de-correo-de-educadores}

\textbf{Asunto}: ``🎯 Nuevo sistema para generar ejercicios matemáticos
únicos''

\textbf{Cuerpo}:

\begin{verbatim}
Hola [Nombre],

¿Cuánto tiempo te toma crear un ejercicio matemático desde cero?

Si eres como la mayoría de educadores, probablemente entre
30 minutos y 2 horas.

¿Y si pudieras generar 300 versiones ÚNICAS del mismo ejercicio
en menos de 5 minutos?

Hoy lanzo el primer demo público de un sistema que hace
exactamente eso: [enlace]

Lo mejor: Es 100% código abierto y gratuito.

Incluye:
→ 5 demos interactivos para probar en tu navegador
→ PDF con 10 versiones para imprimir
→ Código fuente completo
→ Archivo Moodle XML listo para importar

Pruébalo y cuéntame qué te parece.

Saludos,
[Tu nombre]

P.D. Compartiré un nuevo ejercicio cada semana.
¿Qué tema te gustaría ver próximamente?
\end{verbatim}

\begin{center}\rule{0.5\linewidth}{0.5pt}\end{center}

\subsection{🎯 PRÓXIMOS PASOS DESPUÉS DE LA
PUBLICACIÓN}\label{pruxf3ximos-pasos-despuuxe9s-de-la-publicaciuxf3n}

\subsubsection{Semana 1: Engagement
Activo}\label{semana-1-engagement-activo}

\begin{itemize}
\tightlist
\item
  ✅ Responder TODOS los comentarios en \textless{} 2 horas
\item
  ✅ Agradecer cada compartido
\item
  ✅ Conectar con personas que comentaron
\item
  ✅ Unirse a conversaciones relacionadas en otros posts
\end{itemize}

\subsubsection{Semana 2: Análisis y
Optimización}\label{semana-2-anuxe1lisis-y-optimizaciuxf3n}

\begin{itemize}
\tightlist
\item
  ✅ Revisar métricas de LinkedIn Analytics
\item
  ✅ Identificar qué contenido generó más engagement
\item
  ✅ Ajustar estrategia para próximas publicaciones
\item
  ✅ Documentar aprendizajes
\end{itemize}

\subsubsection{Semana 3: Preparar Siguiente
Publicación}\label{semana-3-preparar-siguiente-publicaciuxf3n}

\begin{itemize}
\tightlist
\item
  ✅ Generar nuevo ejercicio
\item
  ✅ Crear demos HTML
\item
  ✅ Preparar contenido visual
\item
  ✅ Escribir post con aprendizajes de la primera publicación
\end{itemize}

\begin{center}\rule{0.5\linewidth}{0.5pt}\end{center}

\subsection{📚 RECURSOS ADICIONALES}\label{recursos-adicionales}

\begin{itemize}
\tightlist
\item
  \textbf{LinkedIn Algorithm Guide 2025}:
  \url{https://www.linkedin.com/pulse/linkedin-2025-content-trends/}
\item
  \textbf{Canva Templates}: Para crear imágenes profesionales
\item
  \textbf{Bitly}: Para acortar y trackear enlaces
\item
  \textbf{Buffer/Hootsuite}: Para programar publicaciones
\end{itemize}

\begin{center}\rule{0.5\linewidth}{0.5pt}\end{center}

\subsection{✅ CHECKLIST FINAL ANTES DE
PUBLICAR}\label{checklist-final-antes-de-publicar}

\begin{verbatim}
□ Post principal escrito y revisado
□ Primer comentario con enlaces preparado
□ Imágenes/carrusel creado (1080x1080 px)
□ Hashtags seleccionados (3-5)
□ Demos HTML funcionando correctamente
□ GitHub Pages activo y accesible
□ Recursos descargables disponibles
□ Horario de publicación programado
□ Notificaciones de LinkedIn activadas
□ Tiempo bloqueado para responder comentarios
\end{verbatim}

\begin{center}\rule{0.5\linewidth}{0.5pt}\end{center}

\textbf{FECHA DE CREACIÓN}: Diciembre 2025 \textbf{BASADO EN}: Mejores
prácticas LinkedIn 2025 \textbf{PRÓXIMA ACTUALIZACIÓN}: Después de
analizar métricas de primera publicación

\end{document}
